\problem{Phonons acoustiques en une dimension}

\question{En utilisant le modèle de Debye, obtenez la dépendance de la chaleur spécifique dans la
température.}

Par définission, la chaleur spécifique est 
\[
    C_V &= \qty(\pdv{E}{T})_V \notag \\
    &= k_B \sum_{k,s} \qty(\frac{\hbar c k}{k_B T})^2 \frac{e^{\beta \hbar c k}}{\qty(e^{\beta \hbar ck} - 1)^2} \notag
\]

Puisque nous somme en 1D et qu'aucuns termes ne dépendent de s, la somme sur s est égale à 1 :
\[
    C_V &=  k_B \sum_{k} \qty(\frac{\hbar c k}{k_B T})^2 \frac{e^{\beta \hbar c k}}{\qty(e^{\beta \hbar ck} - 1)^2} \notag \\
    &= k_B \int_{-k_D}^{k_D} \rho(k) \qty(\frac{\hbar c k}{k_B T})^2 \frac{e^{\beta \hbar c k}}{\qty(e^{\beta \hbar ck} - 1)^2} \, dk \notag \\
    &= \frac{k_B L}{2\pi} \qty(\frac{\hbar c}{k_B T})^2  \int_{-k_D}^{k_D} k^2 \frac{e^{\beta \hbar c k}}{\qty(e^{\beta \hbar ck} - 1)^2} \, dk \notag \\
    &= \frac{k_B L}{2\pi} \qty(\frac{\hbar c}{k_B T})^2  \int_{-k_D}^{k_D} \frac{k^2}{\qty(e^{\beta \hbar ck}e^{\frac{-\beta \hbar ck}{2}} - e^{\frac{-\beta \hbar ck}{2}})^2} \, dk \notag \\
    &= \frac{k_B L}{2\pi} \qty(\frac{\hbar c}{k_B T})^2  \int_{-k_D}^{k_D} \frac{k^2}{\qty(e^{\frac{\beta \hbar ck}{2}} - e^{\frac{-\beta \hbar ck}{2}})^2} \, dk \notag \\
    &= \frac{k_B L}{2\pi} \qty(\frac{\hbar c}{k_B T})^2  \int_{-k_D}^{k_D} \frac{k^2}{4\sinh^2(\frac{\beta \hbar ck}{2})} \, dk \notag
\]

Puisque nous avons une fonction paire, nous pouvons réécrire les bornes de l'intégrale comme 
\[
    C_V = \frac{k_B L}{\pi} \qty(\frac{\hbar c}{k_B T})^2  \int_{0}^{k_D} \frac{k^2}{4\sinh^2(\frac{\beta \hbar ck}{2})} \, dk \notag
\]

Nous pouvons réexprimer le vecteur d'onde de Debye comme $k_D = \frac{T_D}{2T}$ et nous pouvons 
effectuer le changement de variable $x = \frac{\beta \hbar ck}{2}$. Avec cela, nous pouvons obtenir :
\[
    C_V &= \frac{L}{\pi} \frac{\hbar^2 c^2}{k_B T^2} \frac{2 k_B^3 T^3}{\hbar^3c^3} \int_{0}^{\frac{T_D}{2T}} \frac{x^2}{\sinh^2(x)} \, dx \notag \\
    &= \frac{2L}{\pi} \frac{k_B^2 T}{\hbar c} \int_{0}^{\frac{T_D}{2T}} \frac{x^2}{\sinh^2(x)} \, dx \notag
\]

Dans le cas où $T \ll T_D$ :
\[
    C_V =  \frac{2L}{\pi} \frac{k_B^2 T}{\hbar c} \int_{0}^{\infty} \frac{x^2}{\sinh^2(x)} \, dx \notag
\]
Avec Wolfram Alpha, la solution de l'intégrale est $\frac{\pi^2}{6}$. Donc :
\[
    C_V &= \frac{2L}{\pi} \frac{k_B^2 T}{\hbar c} \frac{\pi^2}{6} \notag \\
    &= \frac{L\pi}{3} \frac{k_B^2 T}{\hbar c}
\]

Nous avons bien une solution linéaire en température.
Dans le cas où $T \gg T_D$ nous pouvons développer $\sinh(x) \approx x + \frac{x^3}{6} + ...$. 
En gardant seulement le premier ordre, nous avons :
\[
    C_V &= \frac{2L}{\pi} \frac{k_B^2 T}{\hbar c} \int_{0}^{\frac{T_D}{2T}} \frac{x^2}{x^2} \, dx \notag \\
    &= \frac{2L}{\pi} \frac{k_B^2 T}{\hbar c} \frac{T_D}{2T} \notag  \\
    &= \frac{L k_B^2 T_D}{\pi \hbar c}
\]

Nous avons bien une solution indépendante de la température.

\question{En vous rappelant de la chaleur spécifique
d'un gaz de fermions, indiquez qualitativement comment cela viendra à changer la dépendance
en température $C_V$ , à basse température}

Puisque nous sommes à $T \ll T_F$, la chaleur spécifique pour des fermions est $C_{V_{fermion}} = \frac{Nk_BT}{T_F}$. 
La chaleur spécifique totale est donc :
\[
    C_V &= C_{V_{phonon}} + C_{V_{fermion}} \notag \\
    &= \frac{L\pi}{3} \frac{k_B^2 T}{\hbar c} + \frac{Nk_BT}{T_F} \notag \\
    &= T\qty(\frac{L\pi}{3} \frac{k_B^2}{\hbar c} + \frac{Nk_B}{T_F})
\]

Nous obtenons donc une relation encore linéraire en température, mais avec un facteur dépendent du nombre de fermion.

\clearpage
